\section{Ottimizzazione dell'Analisi Statica: Oracoli per Variabili Non Vincolate}

\subsection{Contesto e abstract compilation}

L'Interpretazione Astratta è una tecnica potente, ma può introdurre
inefficienze significative, specialmente quando si utilizzano domini astratti
relazionali complessi.
Le soluzioni canoniche per migliorare le performance solitamente prevedono
la scelta di un dominio astratto diverso o di operatori astratti differenti.

Tuttavia, esiste un approccio alternativo chiamato \textbf{Abstract
compilation}.

\begin{definizione}{Abstract compilation}{abstract_compilation}

Le tecniche di Abstract compilation consistono nel riscrivere la
rappresentazione intermedia del programma adottata dall'analizzatore.
L'obiettivo è ottimizzare l'esecuzione astratta modificando il codice su cui
l'analisi opererà, cercando un trade-off (compromesso) tra accuratezza ed
efficienza.

\end{definizione}

L'approccio proposto in questa lezione si basa sull'idea che l'efficienza
del processo di valutazione può essere migliorata propagando non solo
l'informazione nota, ma anche l'informazione ``sconosciuta''.

\begin{figure}[H]
    \centering
    \includegraphics[width=0.6\textwidth]{images/th_10/01.png}
    \caption{L'immagine illustra la pipeline dell'analisi statica con
    oracoli per variabili non vincolate.}
    \label{fig:th_10_01}
\end{figure}


\subsection{Likely unconstrained variables}

Il concetto chiave è l'identificazione delle likely unconstrained variables
(LUV), ovvero variabili che molto probabilmente non avranno vincoli
(saranno $\top$, o $[-\infty, +\infty]$) in un determinato punto del programma.

\begin{esempio}{Propagazione dell'incertezza}{luv_propagation}
Considerando l'espressione $x = y + expr$.

Sapendo che $y$ è ``unconstrained'' (cioè $y \to [-\infty, +\infty]$) e
$expr$ ha un valore limitato (es. $[1, 12]$), allora anche il risultato
$x$ sarà necessariamente $[-\infty, +\infty]$.

\begin{figure}[H]
    \centering
    \includegraphics[width=0.8\textwidth]{images/th_10/02.png}
    \caption{L'immagine illustra la propagazione dell'incertezza attraverso
    un'operazione aritmetica.}
    \label{fig:th_10_02}
\end{figure}


Se non si ha alcuna informazione su $y$, è probabile che non si avrà alcuna
informazione sull'intera espressione. Di conseguenza, c'è poco incentivo a
fornire una valutazione astratta precisa (e costosa) per quell'espressione,
poiché il risultato sarà comunque $\top$.
\end{esempio}

L'obiettivo è identificare queste variabili e propagare questa ``mancanza di
informazione'' per semplificare l'analisi.

\subsection{La Pipeline di propagazione LUV}

L'approccio si divide in passaggi distinti, che fungono da pre-analisi
leggera prima dell'analisi target costosa:

\begin{enumerate}
    \item \textbf{Dataflow analysis}: Identificazione delle
    variabili LUV tramite euristiche (Oracoli).
    \item \textbf{Propagation}: Diffusione dell'informazione nel
    CFG (Esistenziale o Universale).
    \item \textbf{Program transformation}: Riscrittura del CFG sostituendo
    le istruzioni complesse con assegnamenti ``havoc''.
\end{enumerate}

\subsubsection{Oracoli e euristiche}

Poiché determinare esattamente se una variabile sarà $\top$ è costoso quanto
l'analisi stessa, si utilizzano degli \textbf{oracoli} euristici.

\begin{figure}[H]
    \centering
    \includegraphics[width=0.5\textwidth]{images/th_10/03.png}
    \caption{L'immagine illustra le diverse tipologie di oracoli.}
    \label{fig:th_10_03}
\end{figure}

Esistono varianti relazionali e non-relazionali.

\begin{definizione}{Oracolo non-relazionale}{non_rel_oracle}
Un oracolo non-relazionale decide se una variabile è LUV basandosi solo sulle
proprietà locali dell'assegnamento.
Ad esempio, per l'operazione $x \leftarrow y \ op \ z$:
\begin{itemize}
    \item Se $y$ è LUV o $z$ è LUV, allora $x$ diventa LUV.
    \item L'oracolo è aggressivo: tende a classificare più variabili come LUV.
\end{itemize}
\end{definizione}

\begin{definizione}{Oracolo relazionale}{rel_oracle}
Un oracolo relazionale cerca di preservare le relazioni tra variabili.
Ad esempio, per $x \leftarrow y \ op \ z$:
\begin{itemize}
    \item Se $op \in \{+, -\}$ e le variabili sono coinvolte, l'oracolo
    potrebbe decidere di non marcare $x$ come LUV immediatamente, per
    preservare potenziali vincoli relazionali.
    \item Questo oracolo è più conservativo (meno falsi positivi, ma meno
    ottimizzazione).
\end{itemize}
\end{definizione}

\begin{nota}{Falsi positivi e negativi}{false_pos_neg}
\begin{itemize}
    \item \textbf{Falso positivo (Perdita di precisione)}: L'oracolo dice
    che $x$ è LUV, ma l'analisi concreta avrebbe trovato un vincolo (es.
    $x \in [0,0]$).
    Questo porta a una perdita di precisione nell'analisi finale.
    \item \textbf{Falso negativo (Mancata ottimizzazione)}: L'oracolo dice
    che $x$ è vincolata, ma in realtà è $[-\infty, +\infty]$.
    L'analisi rimane corretta ma non viene ottimizzata.
\end{itemize}
Indipendentemente dall'euristica, l'analisi finale è sempre \textbf{sound} (corretta).
\end{nota}

\subsubsection{Strategie di Propagazione}

Una volta identificate le LUV localmente, l'informazione deve essere
propagata attraverso i blocchi base del Control Flow Graph (CFG).
I tipi principali di propagazione sono:
\begin{itemize}
    \item \textbf{Propagazione esistenziale ($\exists$)}: Una variabile è nel
    set LUV all'inizio di un blocco se è LUV lungo \emph{almeno uno} degli
    archi entranti.
    \item \textbf{Propagazione universale ($\forall$)}: Una variabile è nel
    set LUV all'inizio di un blocco se è LUV lungo \emph{tutti} gli archi
    entranti.
\end{itemize}

\begin{figure}[H]
    \centering
    \includegraphics[width=0.8\textwidth]{images/th_10/04.png}
    \caption{L'immagine illustra un esempio di propagazione esistenziale e
    universale.}
    \label{fig:th_10_04}
\end{figure}


\subsection{Trasformazione del programma (Havoc)}

L'ultimo passo è la trasformazione del codice intermedio (es. CrabIR).
L'algoritmo prende in input il CFG e la mappa delle variabili LUV post-analisi
($LU_{post}$).
Per ogni istruzione $x \leftarrow expr$:
\begin{itemize}
    \item Se $x \in LU_{post}$ del blocco corrente, l'istruzione viene
    sostituita con $x \leftarrow ?$ (dove $?$ rappresenta un valore non
    deterministico o ``havoc'').
\end{itemize}
Questo rimuove il carico computazionale di valutare $expr$ nell'analisi successiva.


\subsection{Valutazione sperimentale}

L'approccio è stato implementato e testato su vari benchmark.

L'utilizzo di oracoli per variabili non vincolate ha dimostrato:
\begin{itemize}
    \item \textbf{Speedup significativo}: Velocizzazione media
    di 2.61x utilizzando il dominio dei poliedri, con picchi fino a
    8.23x.
    \item \textbf{Precisione}: La perdita di precisione è limitata.
    Con i poliedri si mantiene circa il 90\% dei vincoli originali
    (invarianti).
    \item Le varianti non-relazionali degli oracoli tendono ad essere più
    aggressive e a garantire maggiori miglioramenti di efficienza rispetto
    alla scelta tra propagazione esistenziale o universale.
\end{itemize}
